\section{The finite normal module exists; flatness is a separate gate}
\label{sec:kt-projector-flatness}

The conditional tool KT-NORMAL-FLAT-TRANSFER asks for more
than generic automorphic point detection.  Two of its
integral hypotheses can nevertheless be checked for the
actual source presentation: a finite weight-zero algebra
and a finite normal module surjecting onto the base odd
conormal.  The following statement records exactly this
partial verification, without a flatness assertion.

\begin{lemma}[KT-FINITE-GRADED-NORMAL-MODULE]
Use the weighted presentation of the actual normalized
source ring $S$ in \texttt{ModelOfPres}, together with the
same carried finiteness of \texttt{weightZeroSubalgebra}
used in \texttt{ModelUpper}.  Let $S_j$ be its
weight spaces and let
\[
 A=\bigoplus_{j\in\mathbf Z}S_j\subset S,
 \qquad Q_A=(a-\kappa_S(a):a\in A),
 \qquad B=A/Q_A.
                                                                  \tag{KF.1}
\]
Then $B_0$ is a finite $O$-algebra, and
\[
                 L=(Q_A/Q_A^2)_{-1}                        \tag{KF.2}
\]
is a finite $B_0$-module.  There is a natural surjective
$O$-linear map
\[
                            L\twoheadrightarrow C,         \tag{KF.3}
\]
where $C$ is the odd weight $-1$ cotangent in (KE.2).
The involution acts by minus one on $Q_A/Q_A^2$.
\end{lemma}
\begin{proof}
The source construction supplies a weighted power series
presentation of $S$ with finitely many variables of
weights in $\{-2,-1,0,1,2\}$ and homogeneous kernel.
Precisely, \texttt{ModelOfPres.mop\_cover} constructs the
surjection $q_h:R\to S$ and its finitely generated
homogeneous kernel; \texttt{mop\_main} identifies the
variables with differences of universal matrix entries
and checks the weight bounds.  These are
\texttt{ModelOfPres.lean:34--121,168--238,243--284}.
Composing its weighted presentation of $R$ with $q_h$
gives the asserted presentation of $S$.

The degree-zero image of this presentation is the image
of the ambient model's degree-zero piece.  The complete
proof of
\texttt{ModelIdealFinitenessOfWeightZeroFinite.lean:101--189}
places it inside the already finite
\texttt{weightZeroSubalgebra} by compactness and proves
it finite over $O$.  Hence $S_0$ is finite over $O$.
The standard weighted-presentation argument
\texttt{CommutativeAlgebra/IsWRing.lean:8--56}, theorem
\texttt{isWRing}, makes $A$ a finite-type algebra over
$S_0$.  Its invoked generation theorem is exactly
\texttt{GrGen.lean:72--109}, \texttt{gr\_gen}, whose
homogeneous-series decomposition is proved at lines
15--57: finitely many
nonzero-weight monomials generate all homogeneous
series, with their remaining zero-weight factors in
$S_0$.  Thus $A$ is Noetherian.  The involution commutes
with the torus, so it preserves every $S_j$ and $A$.
Its ideal $Q_A$ is homogeneous and finitely generated
by homogeneous odd elements.  Therefore $B_0$ is the
quotient of $S_0$ by $Q_A\cap S_0$ and is finite over $O$.

The conormal $N_A=Q_A/Q_A^2$ is a finite graded
$B$-module.  Each of its weight pieces is finite over
$B_0$.  Here is a proof that also applies with negative
weights.  For a fixed $j$, form the $B$-submodule of
$N_A$ spanned by $(N_A)_j$.  Since $N_A$ is Noetherian,
finitely many elements $z_1,\ldots,z_r$ of that weight
generate this submodule.  If
$z=\sum_i b_i z_i$ has weight $j$, its weight-$j$
projection is $z=\sum_i(b_i)_0z_i$.  Thus the same
finite list generates $(N_A)_j$ over $B_0$.  This is
the module version of the exact argument in
\texttt{GradedModulesLength/FgHolds.lean:5--45}.
Taking $j=-1$ proves the finiteness of $L$.

The augmentation kills every odd element, so the image
of $Q_A$ lies in $J_0=\ker x_0$ and that of $Q_A^2$
lies in $J_0^2$.  Every element of the odd weight $-1$
piece of $J_0/J_0^2$ has an odd homogeneous lift in
$S_{-1}\subset A$: apply the weight projector and
$(1-\kappa_S)/2$ to any lift.  That lift belongs to
$Q_A$, proving (KF.3).  Finally the proof of
KT-UNSQUARED-NORMAL-JACOBIAN, applied to $A$, shows
that $N_A$ is entirely odd.  This argument needs no
identification of an algebraic conormal with its
completed counterpart beyond the displayed
surjection.
\end{proof}

What is still absent is a proof that $B_0$ is
torsion-free, together with the appropriate reduced
generic structure and full generic normal rank on all
of its components.  The carried rigidity statement
handles the two-sided characteristic-zero orbits; it
does not by itself control vertical contributions or
every fiber of the degree-zero quotient.  Replacing
$B_0$ by its maximal flat reduced quotient may kill
part of $L$ and its image in $C$.  The reverse-map
theorem identifies the latter image with the original
strict Selmer dual, so this loss cannot be ignored.

There are two concrete source boundaries.  First,
\texttt{ModelIdealFinitenessOfWeightZeroFinite} proves
module finiteness only; its displayed conclusions
contain no torsion-freeness or flatness statement.
Second, \texttt{ModelOfPres.mop\_cover} forms the
global ring by imposing the global and marking
relations on a power series extension of the local
ring.  Even the new integral smoothness of the local
ring does not prove that these global relations have
torsion-free quotient.  The characteristic-zero
point detector \texttt{EtHolds} does not fill either
gap, for the reasons in Part I of the projector audit.

\begin{lemma}[KT-FIXED-QUOTIENT-FLATNESS-COUNTEREXAMPLE]
A finite flat reduced augmented $O$-algebra with an
involution can have a fixed-locus quotient with
nonzero $O$-torsion, even when two is a unit.
\end{lemma}
\begin{proof}
For $a=\varpi^n$, $n\geq1$, let $R$ be the subalgebra
of $O^3$ with $O$-basis
\[
  1=(1,1,1),\quad u=(0,a,-a),\quad v=(0,a,a).
                                                                  \tag{KF.4}
\]
These vectors are linearly independent, since $a\ne0$
and two is a unit.  They are closed under
multiplication because
\[
                u^2=av,\qquad uv=au,\qquad v^2=av.          \tag{KF.5}
\]
Thus $R$ is finite free of rank three and reduced,
with generic fiber $K_O^3$.  Projection to its first
coordinate is an augmentation.  The involution
swapping the last two coordinates sends $u$ to
$-u$ and fixes $v$.  Its fixed-locus ideal is $(u)$,
whose $O$-span is generated by $u$ and $av$.
Consequently
\[
                  R/(u)\simeq O[v]/(av,v^2).               \tag{KF.6}
\]
The nonzero class of $v$ is killed by $a$.  In
particular, flatness of an ambient algebra does not
pass to the fixed-locus quotient used in (KF.1).
This is an algebraic scope counterexample, not a
counterexample to the complete source hypotheses.
\end{proof}

No determinant over the actual $B_0$, saturation
identity for $L$, or arithmetic normal Jacobian has
been computed or proved to be a unit.  The verified
finite algebra and module in (KF.1)--(KF.3) identify
where a further integral theorem would apply; they
do not discharge its flatness and support hypotheses.
